\documentclass{article}
\usepackage{amsmath}
\usepackage{amssymb}
\usepackage{graphicx}
\usepackage{hyperref}
\usepackage{enumitem}

\title{DSC 257R: Unsupervised Learning \\ Homework 4}
\author{}
\date{}

\begin{document}

\maketitle

\section*{Informative Projections}

\begin{enumerate}
\item \textbf{Unit vectors.} Find unit vectors in each of the following directions.
    \begin{enumerate}
    \item $\begin{bmatrix}3 \\ 4\end{bmatrix} \in \mathbb{R}^{2}$
    \item $\begin{bmatrix}-4 \\ 12 \\ 3\end{bmatrix} \in \mathbb{R}^{3}$
    \item The all-ones vector in $\mathbb{R}^{d}$
    \end{enumerate}

\item Let $x=\begin{bmatrix}-2 \\ 4\end{bmatrix}$ and $u=\begin{bmatrix}1 \\ 2\end{bmatrix}$.
    \begin{enumerate}
    \item What is the projection of point $x \in \mathbb{R}^{2}$ onto the direction of $u$?
    \item On a grid, plot point $x$, indicate direction $u$, and show the projection.
    \end{enumerate}

\item Let $x=\begin{bmatrix}-4 \\ 3\end{bmatrix} \in \mathbb{R}^{2}$.
    \begin{enumerate}
    \item Find the (unit) direction $u$ such that the projection of $x$ onto $u$ is maximized.
    \item Find the (unit) direction $u$ such that the projection of $x$ onto $u$ is minimized.
    \item Find all (unit) directions $u$ such that the projection of $x$ onto $u$ is zero.
    \end{enumerate}

\item \textbf{The variance in different directions.} A certain data set in $\mathbb{R}^{3}$ has covariance matrix
    \[
    \Sigma = \begin{bmatrix}
    1 & 2 & 0 \\
    2 & 9 & 0 \\
    0 & 0 & 4
    \end{bmatrix}
    \]
    \begin{enumerate}
    \item What is the variance of the data in the direction $\begin{bmatrix}0 \\ 1 \\ 0\end{bmatrix}$?
    \item What is the variance in the direction $\frac{1}{\sqrt{3}}\begin{bmatrix}1 \\ 1 \\ 1\end{bmatrix}$?
    \item The variance in an arbitrary direction $u \in \mathbb{R}^{3}$ is a function of the form
        \[
        a u_{1}^{2} + b u_{2}^{2} + c u_{3}^{2} + d u_{1} u_{2}
        \]
        What is this function exactly? (That is, what are the numerical values of $a, b, c, d$?)
    \end{enumerate}

\item A particular data set in $\mathbb{R}^{5}$ happens to have a diagonal covariance matrix $\Sigma = \operatorname{diag}(5, 10, 20, 1, 4)$. What is the direction of maximum variance?

\item A certain random variable $X \in \mathbb{R}^{3}$ has mean and covariance as follows:
    \[
    \mathbb{E} X = \begin{bmatrix}1 \\ 2 \\ 3\end{bmatrix}, \quad \operatorname{cov}(X) = \begin{bmatrix}
    5 & -3 & 0 \\
    -3 & 5 & 0 \\
    0 & 0 & 4
    \end{bmatrix}
    \]
    \begin{enumerate}
    \item Consider the direction $u = \frac{1}{\sqrt{3}}\begin{bmatrix}1 \\ 1 \\ 1\end{bmatrix}$. What are the mean and variance of $X \cdot u$?
    \item The eigenvectors of $\operatorname{cov}(X)$ can be found in the following list; which ones are they?
        \[
        \begin{bmatrix}1 \\ 0 \\ 0\end{bmatrix}, \begin{bmatrix}0 \\ 1 \\ 0\end{bmatrix}, \begin{bmatrix}0 \\ 0 \\ 1\end{bmatrix}, \quad \frac{1}{\sqrt{2}}\begin{bmatrix}1 \\ 1 \\ 0\end{bmatrix}, \quad \frac{1}{\sqrt{2}}\begin{bmatrix}0 \\ 1 \\ 1\end{bmatrix}, \quad \frac{1}{\sqrt{2}}\begin{bmatrix}1 \\ -1 \\ 0\end{bmatrix}
        \]
    \item Find the eigenvalues corresponding to each of the eigenvectors in part (b). Make it clear which eigenvalue belongs to which eigenvector.
    \item Suppose we used principal component analysis (PCA) to project points $X$ into two dimensions. Which directions would it project onto?
    \item Continuing from part (d), what would be the resulting two-dimensional projection of the point $x = \begin{bmatrix}4 \\ 0 \\ 2\end{bmatrix}$?
    \item Continuing from part (e), suppose that starting from the 2-d projection, we tried to reconstruct the original $x$. What would the three-dimensional reconstruction be, exactly?
    \end{enumerate}

\item \textbf{Some simple visualization methods.} For this problem, we'll again be using the animals with attributes data set (from last week). Recall that this is a small data set that has information about 50 animals. The animals are listed in \texttt{classes.txt}. For each animal, the information consists of values for 85 features: does the animal have a tail, is it slow, does it have tusks, etc. The details of the features are in \texttt{predicates.txt}. The full data consists of a $50 \times 85$ matrix of real values, in \texttt{predicate-matrix-continuous.txt}. Load this real-valued array.
    \begin{enumerate}
    \item We would like to visualize these animals in 2-d. Do this with a PCA projection from $\mathbb{R}^{85}$ to $\mathbb{R}^{2}$. Show the position of each animal, and label each with its name.

        \textbf{Python notes:} You will need to make the plot larger by prefacing your code with
        \begin{verbatim}
from pylab import rcParams
rcParams['figure.figsize'] = 10, 10
        \end{verbatim}
        (or try a different size if this doesn't seem right).

    \item A popular visualization method is the $t$-SNE algorithm. This method takes a numerical parameter called the perplexity and then obtains an embedding by solving a non-convex optimization problem. The t-SNE algorithm is built into scikit-learn. You can invoke it to get a 2-d embedding of a data set $X$ by using:
        \begin{verbatim}
from sklearn.manifold import TSNE
Z = TSNE(n_components=2, perplexity=10.0).fit_transform(X)
        \end{verbatim}
        The perplexity has a significant effect on the output. Try different perplexity values (5, 10, 25, 50) and show each of these embeddings, as you did with the PCA embedding. Bear in mind that t-SNE is a local search algorithm with randomized initialization, and thus even for a fixed perplexity value, different calls to it can return different results.

    \item How can we evaluate these embeddings? Some might seem more visually pleasing than others, or might seem to group animals in a way that agrees more with our own intuitions. Let's look at a somewhat more objective measure. Say we have a data set of $n$ points $x_{1}, \ldots, x_{n} \in \mathbb{R}^{d}$ and we somehow obtain a 2-d visualization $z_{1}, \ldots, z_{n} \in \mathbb{R}^{2}$ of them. How accurately does this 2-d embedding capture the original interpoint distances? To see this,
        \begin{itemize}
        \item Define $D_{i j} = \left\|x_{i} - x_{j}\right\|$ and $\widehat{D}_{i j} = \left\|z_{i} - z_{j}\right\|$.
        \item In general, the $D$ values might be scaled differently from the $\widehat{D}$ values; so define the scaling factor to be $c = \operatorname{mean}(D) / \operatorname{mean}(\widehat{D})$, where $\operatorname{mean}(\cdot)$ denotes the average over all $n^{2}$ entries of the matrix.
        \item The multiplicative factor by which the distance between $x_{i}$ and $x_{j}$ is distorted can be defined by
            \[
            \Delta_{i j} = \max \left(\frac{D_{i j}}{c \cdot \widehat{D}_{i j}}, \frac{c \cdot \widehat{D}_{i j}}{D_{i j}}\right)
            \]
            This ratio is always $\geq 1$.
        \item The average distortion is then $\operatorname{mean}(\Delta)$.
        \end{itemize}
        Compute the average distortion for each of the five embeddings you found (PCA and four t-SNE embeddings). Which of them fares best under this measure?
    \end{enumerate}
\end{enumerate}

\section*{Singular Value Decomposition}

\begin{enumerate}[resume]
\item A particular $4 \times 5$ matrix $M$ has the following singular value decomposition:
    \[
    M = \begin{bmatrix}
    1 & 0 & 0 & 0 \\
    0 & 0 & 0 & 1 \\
    0 & 0 & 1 & 0 \\
    0 & 1 & 0 & 0
    \end{bmatrix}
    \begin{bmatrix}
    4 & 0 & 0 & 0 \\
    0 & 3 & 0 & 0 \\
    0 & 0 & 2 & 0 \\
    0 & 0 & 0 & 1
    \end{bmatrix}
    \begin{bmatrix}
    1 & 0 & 0 & 0 & 0 \\
    0 & 1 & 0 & 0 & 0 \\
    0 & 0 & 1 & 0 & 0 \\
    0 & 0 & 0 & 1 & 0
    \end{bmatrix}
    \]
    Find the best rank-2 approximation to $M$.

\item \textbf{Rank-1 matrices.}
    \begin{enumerate}
    \item Find the best rank-1 approximation to the matrix:
        \[
        M = \begin{bmatrix}
        1 & 2 & 3 \\
        4 & 5 & 6
        \end{bmatrix}
        \]
        If you use the SVD method in Python to solve this problem, you should use the setting \texttt{full\_matrices = 0} to get the kind of decomposition we've been discussing (sometimes called the "thin SVD").
    \item In general, a rank-1 matrix of dimension $p \times q$ is a matrix that can be written in form $u v^{T}$, where $u \in \mathbb{R}^{p}$ and $v \in \mathbb{R}^{q}$. Do you think this decomposition is unique, or is it in general possible to find a different pair of vectors $a \in \mathbb{R}^{p}$ and $b \in \mathbb{R}^{q}$ such that $u v^{T} = a b^{T}$?
    \end{enumerate}
\end{enumerate}

\end{document}
